% BHU PYQ -- Professional Exam Template
\documentclass[11pt,a4paper]{article}

\usepackage[a4paper,top=2.5cm,bottom=2.5cm,left=2.8cm,right=2.8cm,headheight=14pt]{geometry}
\usepackage{amsmath,amssymb}
\usepackage[T1]{fontenc}
\usepackage[utf8]{inputenc}
\usepackage{lmodern}
\usepackage{microtype}
\usepackage{fancyhdr}
\usepackage{enumitem}
\usepackage{listings}
\usepackage{hyperref}
\usepackage{lastpage}
\usepackage{parskip}

\hypersetup{colorlinks=true,urlcolor=black,linkcolor=black,
  pdftitle={BVCA-201: Introduction To Databases -- BHU PYQ}}

\pagestyle{fancy}
\fancyhf{}
\renewcommand{\headrulewidth}{0.4pt}
\renewcommand{\footrulewidth}{0.4pt}
\fancyhead[L]{\small\textit{Banaras Hindu University}}
\fancyhead[R]{\small\textit{BVCA-201: Introduction To Databases}}
\fancyfoot[C]{\small Page \thepage\ of \pageref{LastPage}}

\lstset{
  basicstyle=\small\ttfamily,
  frame=single,
  breaklines=true,
  columns=flexible,
  keepspaces=true
}

\newlist{parts}{enumerate}{2}
\setlist[parts,1]{label=(\alph*),leftmargin=*,itemsep=4pt,topsep=3pt}
\setlist[parts,2]{label=(\roman*),leftmargin=*,itemsep=2pt,topsep=2pt}
\newcommand{\pts}[1]{\hfill\textit{[#1]}}

\begin{document}
\begin{center}
  {\large\bfseries BANARAS HINDU UNIVERSITY}\\[2pt]
  {\normalsize B.Voc. Semester II Examination, 2022-23}\\[6pt]
  \rule{0.8\linewidth}{0.5pt}\\[6pt]
  {\large\bfseries Paper: BVCA-201 --- Introduction To Databases}\\[4pt]
  {\normalsize Time: Three Hours \hfill Maximum Marks: 70}
\end{center}

\rule{\linewidth}{0.4pt}

\smallskip
\noindent\textit{Instructions: Answer the questions in each section as}

\bigskip

| Paper : BVCA-201
Introduction to Databases .
| Note: Answer the questions in each section as
| instructed.
| All question carry marks as indicated.
: Section — A. .
, Note: Long Answer Questions each in 500 words
| approximately (Answer any 2 out of 4. Each
| question carries 12 marks) 2x*12=24
| 1. Draw an E-R Diagram for hospital management
| system in detail and explain.
| 2. What is E-R Diagram ? Explain.
° |

\medskip\noindent\textbf{Question 3.} Whatis normalization ? Explain it with normal |
forms used for normalization. |

\medskip\noindent\textbf{Question 4.} What is the difference between File System
and DBMS ?. Explain any 6 applications of
Database Management System. |

\bigskip\noindent{\large\bfseries SECTION --- B |}
\rule{\linewidth}{0.4pt}\smallskip

Note: Short Answer Questions each in 200 words |
approximately (Answer any 6 out of 9, Each ,
question carries 6 marks) 6x6=36 |

\medskip\noindent\textbf{Question 1.} What is a Database and DBMS ? Explain with
suitable example.

\medskip\noindent\textbf{Question 2.} What is Data Abstraction in DBMS and what
are its three levels ? |

\medskip\noindent\textbf{Question 3.} What are the different types of SOL commands ? |

\medskip\noindent\textbf{Question 4.} What is the difference between TRUNCATE,
DELETE and DROP statements ? .

\medskip\noindent\textbf{Question 5.} What are UNION, MINUS and INTERSECT
commands in DBMS ?

\medskip\noindent\textbf{Question 6.} What is a Trigger in DBMS ?

\medskip\noindent\textbf{Question 7.} What are the advantages of a DBMS ?
2
‘ |

\medskip\noindent\textbf{Question 8.} What is the difference between primary key
and foreign key ?

\medskip\noindent\textbf{Question 9.} What is data model ? List different. types of
‘ data model and explain.
1
, Section-C
Note: Very Short Answer Questions. (Answer all

\medskip\noindent\textbf{Question 10.} Each question carries 1 mark)
1x10=10

\medskip\noindent\textbf{Question 1.} What do you mean by one to many relationships ?
| (a) One class may have many teachers
‘ (b) One teacher can have many classes
\begin{parts}
  \item Many classes may have many teachers
  \item Many teachers may have many classes
\end{parts}

\medskip\noindent\textbf{Question 2.} A Database Management System is a type of
sete eeaveeseeeenee SOftWare.
\begin{parts}
  \item It is a type of system software
  \item It is a kind of application software
  \item It is a kind of general software
  {d) Both (a) and (c) .
\end{parts}

\medskip\noindent\textbf{Question 3.} Which of the following refers to the level lof
data abstraction that describes exactly how
the data actually stored ?
{a} Conceptual Level |
\begin{parts}
  \item Physical Level
  \item File Level
  (d} Logical Level |
\end{parts}

\medskip\noindent\textbf{Question 4.} Which one of the following is a type of Data
Manipulation Command ? |
{a) Create
\begin{parts}
  \item Alter
  \item Delete |
  . (d} All of the above
\end{parts}

\medskip\noindent\textbf{Question 5.} The term “Data independence" refers to |
{a) Data is defined separately and not included
in the programs
\begin{parts}
  \item Programs are not dependent on the logical
  attributes of the data
  \item Programs.are not dependent on the
  physical attributes of the data
  \item Both (b) \& (c)
  4
  . REFNO 065180
\end{parts}

\medskip\noindent\textbf{Question 6.} What is the relation calculus ?
(a) Itis a kind of procedural language
\begin{parts}
  \item Itis a non-procedural language
  \item It is a high-level language
  \_ (d) It is Data Definition language
  ’ 7. Which one of the following refers to the total
  view of the database content ?
  , (a) Conceptual view
(b) Physical view \
  | (c) Internal view
  \item External view
  ) 8. The Database Management Query language is
  generally designed for the ..........s.008. +
  | (a) Support end-users who use English like
  | ; commands
  (b} Specifying the structure of the database
  ! (c}) Support in the development of the complex
  applications software
  -  (d) Allof the above |
  : 5 PTO.
\end{parts}

\medskip\noindent\textbf{Question 9.} Which one of the following is commonly used.
to define the overall design of the database ? }
\begin{parts}
  \item Application program }
  (b} Data definition language
  \item Schema
  \item Source code
  | 10.Which of the following keys is generally used |
  to represents the relationships between the °
  tables ?
  {a} Primary key
  \_ (b) Foreign key . |
  -  {c) Secondary key .
  \item None of the above
  RK
\end{parts}

\vfill
\rule{\linewidth}{0.4pt}
\begin{center}\small
\href{https://bhu-exam.vercel.app/}{Click here to visit our website}\\[4pt]\href{https://me-aryan.vercel.app/}{Click here to visit our contributor}\\[4pt]\href{https://quashan-me.vercel.app/}{Click here to visit our contributor}
\end{center}
\end{document}
